\documentclass[11pt]{article}
\usepackage{amsmath,amssymb}
\usepackage[margin=2.5cm]{geometry}

\newcommand{\ii}{\mathrm{i}}
\newcommand{\ee}{\mathrm{e}}
\newcommand{\orb}{n_{\mathrm{orb}}}

\title{Factored product-grid DFT for the scattering matrix}
\author{}
\date{}

\begin{document}
\maketitle

\section{Definition}

In the M3ph measurement, the scattering matrix $M$ on a 2D target mesh $\{(w^x_t, w^y_t)\}_{t=1}^{N_T}$
is defined by
\begin{equation}\label{eq:M}
  M_t(u,v) \;=\; \sum_{\substack{j:\, u_j^x = u \\ i:\, u_i^y = v}}
    V(j,i)\; \ee^{\ii\, w^x_t\, x_j}\; \ee^{\ii\, w^y_t\, y_i} \,,
\end{equation}
where $j = 0,\dots,k_x{-}1$ and $i = 0,\dots,k_y{-}1$ index the source points
on a product grid with coordinates $(x_j, y_i)$,
orbital labels $(u_j^x, u_i^y) \in \{0,\dots,\orb{-}1\}$,
and $V(j,i)$ is the value matrix (typically $-G^{-1}(j,i)$).

The target frequencies $(w^x_t, w^y_t)$ are fermionic Matsubara frequencies
from a DLR2D or product mesh, with $N_T$ target points.


\section{Naive per-element approach}

The per-element \texttt{buffer\_t<2>} approach computes~\eqref{eq:M}
independently for each orbital pair $(u, v)$.
Buffer $(u, v)$ receives only the subset of source points where
$u_j^x = u$ and $u_i^y = v$, then transforms them to all $N_T$ targets.

The total cost is $\mathcal{O}(N_T \cdot k_x \cdot k_y)$:
for each of the $N_T$ target points, the sum runs over all $k_x \cdot k_y$
source pairs (distributed across $\orb^2$ buffers).
The same $\tau$ values appear in multiple buffers (a given $y_i$ participates
in $k_x$ different pairs), but the per-buffer approach cannot share this work.


\section{Factored approach}

The product-grid structure of the source points allows factoring the 2D
DFT into two 1D operations.

\paragraph{Step 1: Precompute exponential tables.}
Extract the sets of unique target frequencies in each dimension:
$\{w^x_\alpha\}_{\alpha=1}^{N_x}$ and $\{w^y_\beta\}_{\beta=1}^{N_y}$,
with $N_x, N_y \ll N_T$ in general.
For each source point and unique frequency, precompute
\begin{equation}
  E^x(j, \alpha) = \ee^{\ii\, w^x_\alpha\, x_j}\,,\qquad
  E^y(i, \beta)  = \ee^{\ii\, w^y_\beta\,  y_i}\,.
\end{equation}

\paragraph{Step 2: Inner sum (factor out $y$-dimension).}
For each unique $w^y_\beta$, compute the intermediate matrix
\begin{equation}\label{eq:R}
  R_\beta(j, v) \;=\; \sum_{i:\, u_i^y = v} V(j, i)\; E^y(i, \beta)\,.
\end{equation}
This sums over all $k_y$ source points in the $y$-dimension,
scattering by orbital index $v$.
Cost: $\mathcal{O}(N_y \cdot k_x \cdot k_y)$.

\paragraph{Step 3: Outer sum (accumulate over $x$-dimension).}
For each target point $t$ with frequency pair $(w^x_\alpha, w^y_\beta)$,
accumulate
\begin{equation}\label{eq:M_factored}
  M_t(u, v) \;\mathrel{+}=\; \sum_{j:\, u_j^x = u} E^x(j, \alpha)\; R_\beta(j, v)\,.
\end{equation}
This sums over $k_x$ source points per target, scattering by orbital $u$.
Cost: $\mathcal{O}(N_T \cdot k_x \cdot \orb)$.


\section{Complexity comparison}

\begin{center}
\begin{tabular}{lcc}
  \hline
  & Per-element & Factored \\
  \hline
  Exp evaluations & $\mathcal{O}(N_T \cdot (k_x + k_y))$ & $\mathcal{O}(N_x \cdot k_x + N_y \cdot k_y)$ \\
  Multiply-adds   & $\mathcal{O}(N_T \cdot k_x \cdot k_y)$ & $\mathcal{O}(N_y \cdot k_x \cdot k_y + N_T \cdot k_x \cdot \orb)$ \\
  \hline
\end{tabular}
\end{center}

The factored approach wins when $N_y \ll N_T$ and $\orb \ll k_y$, both of which hold
in practice.
For the 4$\times$4 plaquette benchmark ($k = 629$, $N_T = 715$,
$N_x \approx N_y \approx 30$, $\orb = 16$):
factored $\approx 19\text{M}$ ops vs.\ naive $\approx 283\text{M}$ ops ($\sim$15$\times$ reduction).

\end{document}
